% Article VII, additive incorporation, revision 05.
% Authorial architecture: APP--TRIT--TPK memory energy, S01.
% Recovered antecedent: differential and refinement from 30b.
% Collected formalization: note COMPOSICION_VARIACIONAL_TRIT_CONEXION.md.
% Section 30b is retained in full; its anti-Hermitian representative is
% the unitary restriction of the differential developed here.
% This section evaluates the current of a declared action and realization.
% It does not retrospectively select its field or its weight.

\section{Realization of the current in connection coordinates}
\label{vii:sec:composicion-corriente-conexion}

Memory energy assigns each route a response that preserves its
incidences and the order of its transports. Inserting it into a geometric
action requires its components with respect to connection variations.
This step is obtained by differentiating the transport that already
realizes the route: the connection current is the composition of that
differential with the energy covector of the preceding section.

The elliptic, parabolic, and hyperbolic regimes of the TRIT require
retaining the full differential before restricting it to a particular
representation. Unitary rotations constitute one domain of that
construction. The following extension also retains the dilation and
transport directions admitted by the geometric realization being used.

\subsection{Differential on invertible transports}

The fibres, fields, and weight of
\eqref{vii:eq:energia-memoria-discreta} are retained. Now consider
invertible transports and differentiable collections with
\(\dot U_a=A_aU_a\), where \(A_a\) belongs to the tangent space
of the chosen representation. Write
\[
 v_a=U_af_{s(a)},\qquad r_a=f_{t(a)}-v_a,\qquad
 q=\mathsf W_mr,\qquad E_m=\langle r,\mathsf W_mr\rangle.
\]

\begin{proposition}[Full representative of the differential]
\label{vii:prop:diferencial-invertible}
At fixed field and weight,
\begin{equation}
 \mathrm d_UE_m(A)=-2\operatorname{Re}\sum_a
       \langle q_a,A_av_a\rangle
 =\sum_a\operatorname{Re}\operatorname{Tr}(\mathcal G_a^*A_a),
 \qquad \mathcal G_a=-2q_av_a^*.
 \label{vii:eq:gradiente-transporte-completo}
\end{equation}
The anti-Hermitian and Hermitian projections of this representative are
\begin{equation}
 \frac{\mathcal G_a-\mathcal G_a^*}{2}
   =v_aq_a^*-q_av_a^*=\mathcal J_a,
 \qquad
 \frac{\mathcal G_a+\mathcal G_a^*}{2}
   =-(q_av_a^*+v_aq_a^*).
 \label{vii:eq:proyecciones-gradiente-memoria}
\end{equation}
Restriction to anti-Hermitian generators recovers
\eqref{vii:eq:representante-corriente-discreta} exactly.
\end{proposition}

\begin{proof}
The identity \(\dot r_a=-A_av_a\) and self-adjointness of the weight give
\(\dot E_m=2\operatorname{Re}\langle q,\dot r\rangle\).
Furthermore,
\[
 \operatorname{Tr}(\mathcal G_a^*A_a)
 =-2\operatorname{Tr}(v_aq_a^*A_a)=-2q_a^*A_av_a.
\]
The two projections follow by taking the adjoint. For the real
Hilbert--Schmidt inner product, the Hermitian and anti-Hermitian
subspaces are orthogonal; hence the Hermitian part has zero contraction
with an anti-Hermitian generator.
\end{proof}

The formula allows weights that couple edges: their action is computed
in \(q=\mathsf W_mr\) before separating the contributions of each
incidence. If the field or weight varies, the full expression
\eqref{vii:eq:variacion-completa-memoria} is retained; it also holds
for the invertible transports considered here.

\subsection{Quadratic representations of the TRIT}

In a two-dimensional matrix realization of
\(\mathbb A_\tau\), one may take
\[
 J_{+}=\begin{pmatrix}0&1\\-1&0\end{pmatrix},\qquad
 J_{-}=\begin{pmatrix}0&1\\1&0\end{pmatrix},\qquad
 J_{0}=\begin{pmatrix}0&1\\0&0\end{pmatrix}.
\]
Direct multiplication gives
\[
 J_+^2=-I,\qquad J_-^2=I,\qquad J_0^2=0.
\]
Their exponentials realize the elliptic, hyperbolic, and parabolic
regimes, respectively. The first is orthogonal for the fixed positive
inner product. The second preserves the form
\(B=\operatorname{diag}(-1,1)\), since
\(J_-^{\mathsf T}B+BJ_-=0\). The third represents the algebra
of dual numbers. These representations specify examples of the
quadratic types; their use on a screen of another dimension retains
the corresponding representation map.

A calculation distinguishes the components of the differential. On an
edge, take \(U=I\), \(f_s=(1,1)^{\mathsf T}\),
\(f_t=2f_s\), and \(\mathsf W=I\). Then \(q=v=f_s\),
\(\mathcal J=0\), whereas
\begin{equation}
 \mathrm dE(J_-)=-4,\qquad \mathrm dE(J_0)=-2.
 \label{vii:eq:control-regimenes-corriente}
\end{equation}
The full representative retains both responses. This distinction
identifies the space of variations before the matter contraction is
performed.

\subsection{Composition of routes and oriented inversion}

\begin{proposition}[Contragredient transport of the differential]
\label{vii:prop:pullback-invertible}
For \(U_\gamma=U_N\cdots U_1\), let
\(B_k=U_N\cdots U_{k+1}\) and \(B_N=I\). Then
\[
 \dot U_\gamma U_\gamma^{-1}
   =\sum_kB_kA_kB_k^{-1}.
\]
The representative \(\mathcal G_\gamma\) in the final fibre induces,
at incidence \(k\),
\begin{equation}
 \mathcal G_k=B_k^*\mathcal G_\gamma(B_k^{-1})^*.
 \label{vii:eq:transporte-contragrediente}
\end{equation}
\end{proposition}

\begin{proof}
The product rule, followed by multiplication by
\(U_\gamma^{-1}\), proves the first equality. Cyclicity of the
trace gives
\[
 \operatorname{Re}\operatorname{Tr}
  (\mathcal G_\gamma^*B_kA_kB_k^{-1})
 =\operatorname{Re}\operatorname{Tr}
  ((B_k^*\mathcal G_\gamma(B_k^{-1})^*)^*A_k).
\]
For unitary \(B_k\), the conjugation of
\ref{vii:prop:composicion-respuesta} is recovered.
\end{proof}

\begin{proposition}[Inversion by energy congruence]
\label{vii:prop:inversion-congruencia}
For an edge-local energy, inversion is represented by
\begin{equation}
 U_{\bar a}=U_a^{-1},\qquad
 r_{\bar a}=-U_a^{-1}r_a,\qquad
 W_{\bar a}=U_a^*W_aU_a.
 \label{vii:eq:inversion-peso-general}
\end{equation}
It preserves energy. If \(W_a\) remains fixed,
\[
 A_{\bar a}=-U_a^{-1}A_aU_a,\qquad
 \dot W_{\bar a}=U_a^*(A_a^*W_a+W_aA_a)U_a.
\]
The full variation of the inverted energy agrees with the original.
\end{proposition}

\begin{proof}
Substitution into \(r_{\bar a}^*W_{\bar a}r_{\bar a}\)
yields \(r_a^*W_ar_a\). The derivative of the inverse is
\(\dot U_a^{-1}=-U_a^{-1}A_a\); that of the weight follows
from the product rule. Equality of energies holds throughout the
collection, so their derivatives agree, including the term
\(\langle r_{\bar a},\dot W_{\bar a}r_{\bar a}\rangle\).
\end{proof}

Congruence preserves positivity for every invertible transport.
In the unitary representation, it reduces to the inversion rule of
\ref{vii:prop:inversion-respuesta}. The sum uses one orientation per
edge, or half the sum over both orientations.

\subsection{Evaluation of current components}

Let \(\theta=(\theta_b)\) be a real variation of the connection
coordinates of the realization being used. The effective differential
of transport defines the operators
\begin{equation}
 A_a(\theta)=\dot U_aU_a^{-1}
     =\sum_bB_{ab}\theta_b.
 \label{vii:eq:diferencial-conexion-transporte}
\end{equation}
For a finite composition, their contributions are obtained from
\eqref{vii:eq:transporte-contragrediente}. The operators
\(B_{ab}\) are the derivatives of the connection representation
that realizes the route.

\begin{theorem}[Connection current of the memory action]
\label{vii:thm:corriente-conexion-evaluada}
For a matter term \(S_{\rm mem}=\mathfrak a_m E_m\),
where \(\mathfrak a_m\) is its normalization in the action line,
at fixed field, weight, and normalization one has
\begin{equation}
 \delta S_{\rm mem}=-\frac12\sum_b\theta_b s_b,\qquad
 s_b=4\mathfrak a_m\operatorname{Re}
        \sum_a\langle q_a,B_{ab}v_a\rangle.
 \label{vii:eq:componentes-corriente-conexion}
\end{equation}
If the cellular pairing between the connection and current bases is
an invertible matrix \(M\), defined by
\(\delta S_{\rm mem}=-\tfrac12\theta^{\mathsf T}M\sigma\),
the coordinates of the current are
\begin{equation}
 \sigma=M^{-1}s.
 \label{vii:eq:coordenadas-corriente-material}
\end{equation}
\end{theorem}

\begin{proof}
Substitute \eqref{vii:eq:diferencial-conexion-transporte} into
\eqref{vii:eq:gradiente-transporte-completo}, multiply by
\(\mathfrak a_m\), and collect the coefficients of
\(\theta_b\). The variational convention \(-1/2\) fixes the
factor \(4\). Comparing both expressions for every variation
\(\theta\) gives \(M\sigma=s\); invertibility of the pairing
proves existence and uniqueness of these coordinates.
\end{proof}

The formula explicitly realizes the composition mentioned in
\ref{vii:subsec:realizacion-respuesta-discreta}. To evaluate each
component, the field is transported, the residual is computed, the
weight is applied, and contraction is taken with the connection
variation of the same realization. The pairing \(M\) preserves
the orientation and volume of the cells. A diagonal pairing gives the
corresponding division by their volumes; other bases retain the full
matrix.

If the connection also enters \(f\), \(\mathsf W_m\),
or \(\mathfrak a_m\),
\eqref{vii:eq:variacion-completa-memoria} is applied and
\(E_m\delta\mathfrak a_m\) is added. Contributions from all
matter terms are summed before composing the Cartan response.

\begin{corollary}[Compatibility of the current with refinement]
\label{vii:cor:refinamiento-corriente-conexion}
For the compatible collections of
\ref{vii:thm:refinamiento-respuesta}, let \(P_m\) be the fixed
map taking connection variations from level \(m\) to level
\(m+1\). If the memory actions agree on these refined fields and
transports, their covectors satisfy
\[
 s_m=P_m^{\mathsf T}s_{m+1},\qquad
 M_m\sigma_m=P_m^{\mathsf T}M_{m+1}\sigma_{m+1}.
\]
\end{corollary}

\begin{proof}
Differentiating the equality of actions in the direction
\(P_m\theta\) gives
\(-\tfrac12\theta^{\mathsf T}s_m
 =-\tfrac12\theta^{\mathsf T}P_m^{\mathsf T}s_{m+1}\).
The arbitrariness of \(\theta\) and the cellular representations
of both covectors prove the two identities. Composition of these maps
iterates the equality through every finite number of refinements.
\end{proof}

\subsection{Insertion into the Cartan equation}

Section~\ref{vii:sec:corriente-respuesta-cartan} normalizes the
matter three-form by
\(\delta_\omega S_m=-\tfrac12\int\delta\omega^{IJ}\wedge
\sigma_{IJ}\). The preceding theorem provides its cellular
components for the declared memory action and pairing. The Cartan
response uses this current together with the co-tetrad, the
Barbero--Immirzi parameter, and the gravitational coupling of the same
realization.

The dependence of matter on contorsion determines how this equation is
solved. For affine matter, the current is independent of contorsion,
and the linear inversion and quadratic elimination proved in the next
section apply. For a holonomy action, the differential is evaluated at
the current connection, and that dependence is retained when solving
the equation. For example, an edge with
\(f_s=f_t=1\), unit weight, and \(U(t)=e^{it}\) has
\[
 E(t)=2-2\cos t,\qquad E'(t)=2\sin t.
\]
The current here depends on \(t\). The full formula retains this
dependence instead of replacing it by its value at a particular
connection.

Finally, the gravitational density is obtained from the metric
variation of the effective action, with its field and normalization.
The representation of the current, the contorsion it produces, and the
resulting density are the successive operations of this composition.
The subsequent dilution law transports the amplitude thus evaluated
in the corresponding matter section.
